% Generisano: uređuj beleske/sNNN.tex i naslove u manifestu.
\documentclass[11pt,a4paper]{article}
% Zajednički preambulum za dokument za čitanje; nije Beamer tema.
\usepackage[a4paper,margin=20mm,headheight=15pt]{geometry}
\usepackage{lmodern,fontspec}
\usepackage[serbian]{babel}
\setmainfont{Latin Modern Sans}
\setsansfont{Latin Modern Sans}
\setmonofont{Latin Modern Mono}
\usepackage{amsmath,amssymb,mathtools,graphicx,xcolor}
\usepackage{array,booktabs,tabularx,multirow,makecell,listings,fancyhdr}
\pagestyle{fancy}\fancyhf{}
\newcommand{\BeleskeTekuciID}{NEOZNACEN}
\fancyhead[L]{\small Beleške uz slajd \texttt{\expandafter\detokenize\expandafter{\BeleskeTekuciID}}}
\fancyfoot[R]{\small\thepage}
\setlength{\parindent}{0pt}\setlength{\parskip}{6pt}
\newwrite\BeleskeMapa
\immediate\openout\BeleskeMapa=\jobname.notes-map.tsv
\AddToHook{shipout/before}{%
  \immediate\write\BeleskeMapa{\BeleskeTekuciID\space\thepage}}
\AddToHook{enddocument/afterlastpage}{\immediate\closeout\BeleskeMapa}
\newcommand{\BeleskeID}[1]{%
  \def\BeleskeZadatiID{#1}%
  \ifx\BeleskeTekuciID\BeleskeZadatiID\else
    \PackageError{beleske}{Pogresan ID beleske: #1}{Proveri mapiranje slajda.}%
  \fi}
\newcommand{\BeleskaSlajda}[4]{%
  \clearpage\gdef\BeleskeTekuciID{#1}%
  {\Large\bfseries Slajd \textnormal{\texttt{\detokenize{#1}}}: #3\par}\medskip
  \includegraphics[page=#2,width=\linewidth]{\BeleskePrezentacija}\par\medskip
  \input{#4}}

\newcommand{\BeleskePrezentacija}{build/09_3_mos_logicka_kola.pdf}
\begin{document}
\BeleskaSlajda{09_3-s001}{1}{Logička kola sa MOS tranzistorima — 3. deo}{beleske/s001.tex}
\BeleskaSlajda{09_3-s002}{2}{Dinamičke karakteristike CMOS invertora}{beleske/s002.tex}
\BeleskaSlajda{09_3-s003}{3}{Kapacitivnosti MOS tranzistora}{beleske/s003.tex}
\BeleskaSlajda{09_3-s004}{4}{Kapacitivnost gejta i režim rada}{beleske/s004.tex}
\BeleskaSlajda{09_3-s005}{5}{Dva problema optimizacije}{beleske/s005.tex}
\BeleskaSlajda{09_3-s006}{6}{Optimizacija kašnjenja CMOS invertora}{beleske/s006.tex}
\BeleskaSlajda{09_3-s007}{7}{Kapacitivnosti u lancu invertora}{beleske/s007.tex}
\BeleskaSlajda{09_3-s008}{8}{Uticaj parazitnih kapacitivnosti}{beleske/s008.tex}
\BeleskaSlajda{09_3-s009}{9}{Srednje kašnjenje i ekvivalentne otpornosti}{beleske/s009.tex}
\BeleskaSlajda{09_3-s010}{10}{Milerova kapacitivnost i skaliranje}{beleske/s010.tex}
\BeleskaSlajda{09_3-s011}{11}{Kašnjenje u zavisnosti od širine kanala}{beleske/s011.tex}
\BeleskaSlajda{09_3-s012}{12}{Optimalni odnos širina kanala}{beleske/s012.tex}
\BeleskaSlajda{09_3-s013}{13}{Jedinični CMOS invertor}{beleske/s013.tex}
\BeleskaSlajda{09_3-s014}{14}{Optimizacija lanca invertora}{beleske/s014.tex}
\BeleskaSlajda{09_3-s015}{15}{Lanac skaliranih invertora}{beleske/s015.tex}
\BeleskaSlajda{09_3-s016}{16}{Kašnjenje prvog stepena}{beleske/s016.tex}
\BeleskaSlajda{09_3-s017}{17}{Efektivni fanout}{beleske/s017.tex}
\BeleskaSlajda{09_3-s018}{18}{Optimalno geometrijsko skaliranje}{beleske/s018.tex}
\BeleskaSlajda{09_3-s019}{19}{Fiksiran broj stepena}{beleske/s019.tex}
\BeleskaSlajda{09_3-s020}{20}{Optimalni broj stepena}{beleske/s020.tex}
\BeleskaSlajda{09_3-s021}{21}{Broj stepena za faktor 3,6}{beleske/s021.tex}
\BeleskaSlajda{09_3-s022}{22}{Faktor povećanja četiri}{beleske/s022.tex}
\BeleskaSlajda{09_3-s023}{23}{Nebaferisana i dvostruko baferisana kola}{beleske/s023.tex}
\BeleskaSlajda{09_3-s024}{24}{Povećanje kapaciteta i rasterećenje}{beleske/s024.tex}
\BeleskaSlajda{09_3-s025}{25}{Složena CMOS logička kola}{beleske/s025.tex}
\BeleskaSlajda{09_3-s026}{26}{Realizacije PUN i PDN mreža}{beleske/s026.tex}
\BeleskaSlajda{09_3-s027}{27}{Dimenzije dvoulaznih NI i NILI kola}{beleske/s027.tex}
\BeleskaSlajda{09_3-s028}{28}{Dimenzionisanje NILI kola}{beleske/s028.tex}
\BeleskaSlajda{09_3-s029}{29}{Dimenzionisanje NI kola}{beleske/s029.tex}
\BeleskaSlajda{09_3-s030}{30}{Poređenje površina NI i NILI kola}{beleske/s030.tex}
\BeleskaSlajda{09_3-s031}{31}{Sinteza PDN mreže}{beleske/s031.tex}
\BeleskaSlajda{09_3-s032}{32}{Sinteza PUN mreže}{beleske/s032.tex}
\BeleskaSlajda{09_3-s033}{33}{Dualnost PUN i PDN mreža}{beleske/s033.tex}
\BeleskaSlajda{09_3-s034}{34}{Složeno CMOS kolo sa dimenzijama}{beleske/s034.tex}
\BeleskaSlajda{09_3-s035}{35}{Logički trud i kritična putanja}{beleske/s035.tex}
\BeleskaSlajda{09_3-s036}{36}{Karakteristična vremenska konstanta}{beleske/s036.tex}
\BeleskaSlajda{09_3-s037}{37}{Normalizovano kašnjenje lanca}{beleske/s037.tex}
\BeleskaSlajda{09_3-s038}{38}{Vremenska konstanta NI kola}{beleske/s038.tex}
\BeleskaSlajda{09_3-s039}{39}{Vremenska konstanta NILI kola}{beleske/s039.tex}
\BeleskaSlajda{09_3-s040}{40}{Mešovita kritična putanja}{beleske/s040.tex}
\BeleskaSlajda{09_3-s041}{41}{Optimizacija dva susedna kola}{beleske/s041.tex}
\BeleskaSlajda{09_3-s042}{42}{Jednaki proizvodi vremena i fanouta}{beleske/s042.tex}
\BeleskaSlajda{09_3-s043}{43}{Definicija logičkog truda}{beleske/s043.tex}
\BeleskaSlajda{09_3-s044}{44}{Opšti izraz logičkog truda}{beleske/s044.tex}
\BeleskaSlajda{09_3-s045}{45}{Zašto ne spajamo obične izlaze}{beleske/s045.tex}
\BeleskaSlajda{09_3-s046}{46}{CMOS izlaz sa tri stanja}{beleske/s046.tex}
\BeleskaSlajda{09_3-s047}{47}{NMOS izlaz sa otvorenim drejnom}{beleske/s047.tex}
\BeleskaSlajda{09_3-s048}{48}{Zajednička linija više predajnika}{beleske/s048.tex}
\BeleskaSlajda{09_3-s049}{49}{Visoka impedansa i dodatna kašnjenja}{beleske/s049.tex}
\BeleskaSlajda{09_3-s050}{50}{PTL — prolazna logička kola}{beleske/s050.tex}
\BeleskaSlajda{09_3-s051}{51}{Prenos jedinice i nule kroz NMOS}{beleske/s051.tex}
\BeleskaSlajda{09_3-s052}{52}{Selektorska prolazna logika}{beleske/s052.tex}
\BeleskaSlajda{09_3-s053}{53}{Prolazna I i ILI kola}{beleske/s053.tex}
\BeleskaSlajda{09_3-s054}{54}{Degradacija naponskih nivoa}{beleske/s054.tex}
\BeleskaSlajda{09_3-s055}{55}{Pravilna veza prolaznih tranzistora}{beleske/s055.tex}
\BeleskaSlajda{09_3-s056}{56}{Restauracija naponskih nivoa}{beleske/s056.tex}
\BeleskaSlajda{09_3-s057}{57}{Multiplekser u prolaznoj logici}{beleske/s057.tex}
\BeleskaSlajda{09_3-s058}{58}{Pomerač — množenje i deljenje sa 2}{beleske/s058.tex}
\BeleskaSlajda{09_3-s059}{59}{Barrel shifter — množenje, deljenje sa 2 na n-ti stepen}{beleske/s059.tex}
\BeleskaSlajda{09_3-s060}{60}{Transmisioni gejt}{beleske/s060.tex}
\BeleskaSlajda{09_3-s061}{61}{Bilateralni prekidač: negativan ulaz}{beleske/s061.tex}
\BeleskaSlajda{09_3-s062}{62}{Bilateralni prekidač: pozitivan ulaz}{beleske/s062.tex}
\BeleskaSlajda{09_3-s063}{63}{Bilateralni prekidač: ulaz na nuli}{beleske/s063.tex}
\BeleskaSlajda{09_3-s064}{64}{Otpornost transmisionog gejta}{beleske/s064.tex}
\BeleskaSlajda{09_3-s065}{65}{Isključen bilateralni prekidač}{beleske/s065.tex}
\BeleskaSlajda{09_3-s066}{66}{Transmisioni gejt u digitalnoj logici}{beleske/s066.tex}
\BeleskaSlajda{09_3-s067}{67}{Koju funkciju realizuje kolo?}{beleske/s067.tex}
\BeleskaSlajda{09_3-s068}{68}{Dinamička logika}{beleske/s068.tex}
\BeleskaSlajda{09_3-s069}{69}{Faze rada dinamičke logike}{beleske/s069.tex}
\BeleskaSlajda{09_3-s070}{70}{Problem kaskade dinamičkih kola}{beleske/s070.tex}
\BeleskaSlajda{09_3-s071}{71}{Domino logika}{beleske/s071.tex}
\end{document}
